\documentclass[12pt,a4paper]{article}
\usepackage{geometry}
\geometry{left=2.5cm,right=2.5cm,top=2.0cm,bottom=2.5cm}
\usepackage[english]{babel}
\usepackage{amsmath,amsthm}
\usepackage{amsfonts}
\usepackage[longend,ruled,linesnumbered]{algorithm2e}
\usepackage{fancyhdr}
\usepackage{ctex}
\usepackage{array}
\usepackage{listings}
\usepackage{color}
\usepackage{graphicx}

\begin{document}


\title{
{《优化理论与算法》第$6$次作业
}
}
\date{}

% \author{
% 姓名：\underline{你的名字}~~~~~~
% 学号：\underline{你的学号}~~~~~~}

\maketitle
\begin{itemize}
	\item 作业截止于{\bf 周五（4/17/2026）}，在Canvas系统中提交PDF或照片文件
	\item 作业题目的tex文件可以在Canvas中下载
	\item 鼓励相互讨论，但不要抄袭.
\end{itemize}
%%%%%%%%%%%%%%%%%%%%%%%%%%%%%%%%%%%%%%%%%%%%%%%%%%%%%%%%%%%%%%%%%
%%%%%%%%%%%%%%%%%%%%%%%%%%%%%%%%%%%%%%%%%%%%%%%%%%%%%%%%%%%%%%%%%


\section{习题}


\paragraph{2.1}
给定一个有向网络 $G=(V,E)$，其中顶点集合 $V$ 包含 $6$ 个节点，分别记为
\[
\{s,\,1,\,2,\,3,\,4,\,t\},
\]
其中 $s$ 为源点，$t$ 为汇点。下表给出了网络中各条有向边及其容量：

\[
\begin{array}{c|cc}
\textbf{边} & \textbf{容量} \\ \hline
(s,1) & 15 \\
(s,2) & 5 \\
(1,2) & 5 \\
(1,3) & 10 \\
(2,4) & 8 \\
(3,4) & 7 \\
(3,t) & 6 \\
(4,t) & 12
\end{array}
\]

请根据以上数据完成下列要求：

\begin{enumerate}
  \item \textbf{模型建立}：  
  建立最大流问题模型，建立最小割问题模型

  \item \textbf{求解}：  
  使用任何合适的求解器，求解上述最大流模型，请给出最优解中各条边的流量分配结果；求解上述最小割模型，请给出最优分割。最后，验证最大流最小割定理成立。
  

  \item \textbf{结果分析与讨论}：
  \begin{enumerate}
    \item 若将边 $(1,3)$ 的容量由 10 调整为 5，最大流量是否会发生变化？若会，新的最优流量值是多少？  
    \item 若在节点 $2$ 与节点 $3$ 之间增设一条新边 $(2,3)$，容量为 6，则新的最优流量是多少？  
  \end{enumerate}
\end{enumerate}


\paragraph{2.2} 
试证明：引入容量下界的最小成本流问题，也符合整数定理，即：
如果如下线性规划中，
\begin{equation*}
\begin{aligned}
& {\text{minimize}}
& & \sum_{(i,j) \in E} c_{ij}x_{ij} \\
& \text{subject to}
& & \sum_{(k,i) \in E} x_{ki} - \sum_{(i,j)\in E} x_{ij} =d_i,&\forall i \in V\\
& & & l_{ij} \leq x_{ij} \leq u_{ij} ,&\forall (i,j) \in E\\
\end{aligned}
\end{equation*}
参数$d_i, l_{ij}, u_{ij}$均为整数，那么其所有基本可行解均为整数解。

\end{document}